% Integration-ready: amsmath, amssymb, amsthm, booktabs, hyperref.
% No priority claim beyond the pinned repository and inspected incoming files.
\section{Exact sampling of the unresolved mixed Gaussian generator}
\label{mix:sec}

The mixed Gaussian sums
\[
 S_p=\sum_{n=0}^{\infty}\frac{(-1)^nH_n}{(2n+1)^p}
\]
have a proved reduction at $p=2,4$, while the fixed-weight reductions at
$p=6,8$ remain conjectural in the pinned manuscript.  The incoming
uniform-continuation report additionally supplies conjectures at $p=10,12$.
The following results do not settle these fixed-weight questions. They give
an exact arithmetic sampling theorem for the whole generator, explicit
finite parts at its poles, and convergent identities involving the even
indices. A new, separately validated candidate at $p=14$ follows below.

Write $L=\log2$, $G=\beta(2)$ and
\[
 \mathcal S(t)=\sum_{p\ge1}S_pt^{p-1},\qquad
 \mathcal O(t)=\frac{\mathcal S(t)-\mathcal S(-t)}2
             =\sum_{m\ge1}S_{2m}t^{2m-1}.
\]
The canonical manuscript already proves the meromorphic continuation of
$\mathcal S$, its Mellin integral and its simple poles. We retain these
facts as prerequisites, and include the short derivations needed for the
new formulas. In particular, no new assertion of period independence or
of minimal depth is made.

\subsection{A shift equation and rational sampling with all parameter jets}

Put
\[
 B(t)=\int_0^1\frac{x^{-t}}{1+x^2}\,dx
 =\frac14\left\{\psi\left(\frac{3-t}{4}\right)
                  -\psi\left(\frac{1-t}{4}\right)\right\},
 \qquad \Re t<1,
\]
and continue $B$ meromorphically. Its elementary recurrence and reflection
formula are
\begin{equation}\label{mix:eq:B}
 B(t)+B(t-2)=\frac1{1-t},\qquad
 B(t)+B(-t)=\frac{\pi}{2\cos(\pi t/2)}.
\end{equation}
The first follows by adding the two integrands; the second follows by
splitting $\int_0^\infty x^{-t}/(1+x^2)\,dx$ at one.

\begin{theorem}[Meromorphic shift equation]\label{mix:thm:shift}
On the complex plane, as an identity of meromorphic functions,
\begin{equation}\label{mix:eq:shift}
 \boxed{\quad
 \mathcal S(t)+\mathcal S(t-2)
       =\frac{2B(t-2)-L}{1-t}.
 \quad}
\end{equation}
Removable values and finite parts are taken after combining the terms.
\end{theorem}
\begin{proof}
Let $q(x)=\log(1+x^2)/(1+x^2)$. The defining generating integral is
$\mathcal S(t)=-\int_0^1x^{-t}q(x)\,dx$ for $\Re t<3$.
Adding its shift by two gives
$-\int_0^1x^{-t}\log(1+x^2)\,dx$.
For $t\ne1$, integration by parts gives precisely
\eqref{mix:eq:shift}; the lower boundary term vanishes because
$x^{1-t}\log(1+x^2)=O(x^{3-\Re t})$.
The identity first holds in an open half-plane and extends by the
meromorphic identity theorem. The point $t=1$ is removable on both sides.
\end{proof}

We use the series convention
\[
 \operatorname{Li}_{a,b}(x,y)
 =\sum_{n>m\ge1}\frac{x^ny^m}{n^am^b}.
\]
For $d\ge1$ let $\mathcal R_d=\{\xi:\xi^d=-1\}$.
No member of $\mathcal R_d$ equals one. At unit-modulus arguments,
we use the radial boundary values of the multiple polylogarithms; all
values below also have ordinarily convergent defining series.
Indeed, for $r\ge1$ the double series is absolutely convergent.
For $r=0$, write its inner argument as $y=\xi/\eta$ and outer
argument as $x=\xi^{-1}\ne1$. If $y\ne1$, the inner sum is
$\Li_1(y)+O(n^{-1})$, so Dirichlet's test treats the main outer
series and the remainder converges absolutely. If $y=1$,
$H_{n-1}/n$ decreases eventually to zero, and the same test applies.

\begin{theorem}[Cyclotomic rational sampling and parameter jets]
\label{mix:thm:rational}
Let $1\le A\le d$ be integers, put $\alpha=1-2A/d$, and let $r\ge0$
be an integer. Then
\begin{align}
 \frac{\mathcal S^{(r)}(\alpha)}{r!}
 &=\frac12\left(\frac d2\right)^r
   \sum_{\xi,\eta\in\mathcal R_d}
      \xi^A\operatorname{Li}_{r+1,1}
                    \left(\xi^{-1},\frac\xi\eta\right).
 \label{mix:eq:rational}\\[-1mm]
 \frac{\mathcal S^{(r)}(\alpha)}{r!}
 &=\sum_{k\ge0}\binom{k+r}{r}S_{k+r+1}\alpha^k.
 \label{mix:eq:allweight}
\end{align}
Thus every rational sample in $[-1,1)$ and every parameter jet is an
explicit $\mathbb Q(\zeta_{2d})$-linear combination of cyclotomic doubles
of the fixed weight $r+2$.  All sums on the right of
\eqref{mix:eq:rational} are finite; complex-conjugate terms make the
result real. Choosing $A/d=(1-\alpha)/2$ in lowest terms gives level $2d$
as an upper bound, not a claim of minimal conductor.
The shift equation extends these evaluations to every rational sample
at which $\mathcal S$ is regular, and gives Laurent coefficients at its
remaining rational points.
\end{theorem}
\begin{proof}
Substitute $u=x^2$ in the Mellin integral and then $u=v^d$. This gives
\[
 \mathcal S(\alpha)
 =-\frac d2\int_0^1
       \frac{v^{A-1}\log(1+v^d)}{1+v^d}\,dv.
\]
The exact partial fraction identity is
\begin{equation}\label{mix:eq:partial}
 \frac{d\,v^{A-1}}{1+v^d}
       =-\sum_{\xi\in\mathcal R_d}\frac{\xi^A}{v-\xi}.
\end{equation}
Both sides are proper rational functions with the same simple poles and
residues: the residue on the left at $\xi$ is
$\xi^{A-d}=-\xi^A$. Their difference therefore vanishes.
Along $0\le v\le1$,
\[
 \log(1+v^d)=\sum_{\eta\in\mathcal R_d}\log(1-v/\eta).
\]
The logarithms are their principal values, continued from $v=0$; equality
follows by differentiating and matching the value at zero. None of the
factors meets its branch cut on this path.

Differentiating $r$ times with respect to $t$ inserts
$(-\log x)^r=(d/2)^r(-\log v)^r$. Finally, the standard convergent
iterated-integral formula, with the displayed series convention, is
\[
 \operatorname{Li}_{r+1,1}(\xi^{-1},\xi/\eta)
 =\frac1{r!}\int_0^1
       \frac{(-\log v)^r\log(1-v/\eta)}{v-\xi}\,dv.
\]
Substitution proves \eqref{mix:eq:rational}. All integrals converge
absolutely; at zero the logarithm is $O(v)$, and the root denominators
do not vanish on $[0,1]$. Parameter differentiation is justified on
compact subsets of $\Re t<3$ by integrable powers of $|\log x|$.
Taylor differentiation inside $|t|<3$ proves
\eqref{mix:eq:allweight}.  Repeated application and differentiation of
\eqref{mix:eq:shift} proves the last assertion.
\end{proof}

For example, at $\alpha=1/3$ one may take $(A,d)=(1,3)$, so the
whole infinite sum $\sum_{p\ge1}S_p3^{1-p}$ is a finite combination
of weight-two values at sixth roots of unity. At $\alpha=1/2$, one takes
$(A,d)=(1,4)$ and obtains level eight. This is a fixed-weight evaluation
of an infinite generating sum; it is not a termwise reduction of the
individual periods $S_p$.

\subsection{Elementary values and finite parts of the odd generator}

Define the rational finite sums
\[
 a_N=\sum_{j=1}^N\frac{(-1)^{j-1}}{2j-1},\quad
 b_N=\sum_{j=1}^N\frac{a_{j-1}}{2j-1},\quad
 c_N=\sum_{j=1}^N\frac{(-1)^{j-1}}{(2j-1)^2},
\]
with empty sums zero. Also put
\[
 \bar H_N^{(s)}=\sum_{j=1}^N\frac{(-1)^{j-1}}{j^s},\qquad
 e_N=\sum_{j=1}^N\frac{\bar H_{j-1}^{(1)}}j,
 \qquad P=\frac{L^2}{4}-\frac{\pi^2}{48}.
\]
The notation $\operatorname{FP}_{t=t_0}$ denotes the constant term in
the Laurent expansion in the local coordinate $t-t_0$.

\begin{theorem}[All integer samples of the unresolved parity component]
\label{mix:thm:integer}
For every integer $N\ge0$,
\begin{equation}\label{mix:eq:even-sample}
 \boxed{\quad
 \mathcal O(2N)=(-1)^N\{c_N+2b_N-a_NL\}.
 \quad}
\end{equation}
At $t=2N+1$, its residue is
$(-1)^{N+1}H_N/2$, and its finite part is
\begin{align}
 \operatorname{FP}_{t=2N+1}\mathcal O(t)
 ={}&(-1)^N\left\{
 P-\frac L2(H_N+\bar H_N^{(1)})\right.\notag\\
 &\left.\hspace{13mm}
 +\frac{H_N^{(2)}+\bar H_N^{(2)}}4+\frac{e_N}2
 \right\}.
 \label{mix:eq:odd-fp}
\end{align}
For $N=0$ the pole is absent, and \eqref{mix:eq:odd-fp} is the
ordinary value $\mathcal O(1)=P$.
\end{theorem}
\begin{proof}
First, the elementary beta/reflection evaluation gives
$\mathcal S(0)=G-\pi L/2$. Direct integration after $u=x^2$ gives
\[
 \mathcal S(-1)=-\frac{L^2}4,\qquad
 \mathcal S(1)=\frac{L^2}4-\frac{\pi^2}{24}.
\]
The recurrence \eqref{mix:eq:B}, starting with $B(0)=\pi/4$, yields
\[
 B(2j)=(-1)^j(\pi/4+a_j),\qquad
 B(-2j)=(-1)^j(\pi/4-a_j).
\]
Let $h_N^{\mathrm o}=\sum_{j=1}^N(2j-1)^{-1}$. Iterating
\eqref{mix:eq:shift} at positive and negative even integers gives
\begin{align*}
 \mathcal S(2N)&=(-1)^N\{
 \mathcal S(0)-La_N+(\pi/2)h_N^{\mathrm o}+2b_N\},\\
 \mathcal S(-2N)&=(-1)^N\{
 \mathcal S(0)+La_N+(\pi/2)h_N^{\mathrm o}-2b_N-2c_N\}.
\end{align*}
Subtracting proves \eqref{mix:eq:even-sample}.

For the odd finite parts, the normally convergent pole expansion of $B$
or its digamma formula gives, for $j\ge0$,
\[
 \operatorname*{Res}_{t=2j+1}B(t)=(-1)^{j+1},\qquad
 \operatorname{FP}_{t=2j+1}B(t)
       =\frac{(-1)^j}{2}(\bar H_j^{(1)}-L).
\]
Write $\mathcal C_N^{\mathrm{fp}}=\operatorname{FP}_{t=2N+1}\mathcal S(t)$.
Extracting the
constant term of \eqref{mix:eq:shift} gives, for $N\ge1$,
\[
 \mathcal C_N^{\mathrm{fp}}+\mathcal C_{N-1}^{\mathrm{fp}}
 =\frac{(-1)^N}{2N^2}
  +\frac{(-1)^N\bar H_{N-1}^{(1)}}{2N}
  +\frac{1-(-1)^N}{2N}L.
\]
The first term is essential: it comes from the pole of $B(t-2)$ times
the linear term of $(1-t)^{-1}$.
Starting with $\mathcal C_0^{\mathrm{fp}}=L^2/4-\pi^2/24$,
the recurrence gives
\[
 \mathcal C_N^{\mathrm{fp}}=(-1)^N\left\{
 \mathcal C_0^{\mathrm{fp}}+\frac{H_N^{(2)}}2+\frac{e_N}2
       -\frac L2(H_N+\bar H_N^{(1)})\right\}.
\]
On the negative odd axis, polynomial division of $u^N/(1+u)$ gives
\[
 \mathcal S(-(2N+1))=(-1)^N\left\{
 -\frac{L^2}4+\frac L2(H_N+\bar H_N^{(1)})
       -\frac{e_N+\bar H_N^{(2)}}2\right\}.
\]
Taking half the difference proves \eqref{mix:eq:odd-fp}.
The residue is half the already known residue of $\mathcal S$ at the
positive pole; $\mathcal S(-t)$ is regular there.
\end{proof}

\subsection{Convergent identities after subtracting known poles}

\begin{corollary}[Three explicit identities across all even indices]
\label{mix:cor:resummation}
The following series converge absolutely and satisfy
\begin{align}
 \sum_{m\ge1}S_{2m}
 &=\frac{L^2}4-\frac{\pi^2}{48},\label{mix:eq:sum1}\\
 \sum_{m\ge1}2^{2m-1}S_{2m}
 &=L-1,\label{mix:eq:sum2}\\
 \sum_{m\ge1}3^{2m-1}(S_{2m}+3^{-2m})
 &=L-\frac{L^2}4+\frac{\pi^2}{48}-\frac7{12}.
 \label{mix:eq:sum3}
\end{align}
\end{corollary}
\begin{proof}
The first two are $\mathcal O(1)$ and $\mathcal O(2)$, inside the
Taylor disk of radius three. The $n=1$ term of the defining harmonic
sum contributes $t/(t^2-9)$ to $\mathcal O(t)$. Subtracting it leaves
a function analytic in $|t|<5$. Its value at $t=3$ is
\[
 \operatorname{FP}_{t=3}\mathcal O(t)-\frac1{12}
 =-P+L-\frac12-\frac1{12},
\]
which is the last right side. This also explains why adding $3^{-2m}$
before summation is necessary at this boundary point.
\end{proof}

More generally, for an integer $K\ge0$, set
\[
 T_p^{[K]}=S_p-\sum_{n=1}^K\frac{(-1)^nH_n}{(2n+1)^p},\qquad
 \mathcal O^{[K]}(t)=\mathcal O(t)
   -\sum_{n=1}^K\frac{(-1)^nH_nt}{(2n+1)^2-t^2}.
\]
Then $\mathcal O^{[K]}(t)=\sum_{m\ge1}T_{2m}^{[K]}t^{2m-1}$ has
Taylor radius exactly $2K+3$. Theorem~\ref{mix:thm:integer} supplies
its value at $t=2K+2$ and its regular value at $t=2K+1$ after the
indicated finite rational pole contributions are removed. This gives an
infinite family of convergent identities with elementary right sides.
For $d=2K+3$, $|t|<d$, and $M\ge0$, a useful explicit remainder bound is
\begin{equation}\label{mix:eq:tail}
 \left|\sum_{m>M}T_{2m}^{[K]}t^{2m-1}\right|
 \le \frac{H_{K+1}|t|^{2M+1}}
 {d^{2M+2}(1-|t|^2/d^2)}.
\end{equation}
Indeed $H_n/(2n+1)^p$ strictly decreases for $n\ge1,p\ge2$, so the
alternating-series estimate gives $|T_p^{[K]}|\le H_{K+1}/d^p$.
Summing the resulting geometric series proves the bound. The actual
nonzero poles at $\pm d$ prove the exact Taylor radius.

These formulas impose exact analytic consistency constraints on an
entire proposed sequence of fixed-weight reductions. They do not prove
any individual $S_{2m}$ reduction, because the resummations combine
different weights. The distinction matters when they are used to audit
conjectural coefficient families.

\section{A new frozen mixed Gaussian candidate at weight fifteen}
\label{mix:sec:S14}

Put $g_{a,b}=\operatorname{Im}\operatorname{Li}_{a,b}(i,1)$.
The inspected incoming reports reach $S_{12}$; the present search advances
this particular candidate family to $S_{14}$.  The old $S_6$ and $S_8$
relations, and the incoming $S_{10}$ and $S_{12}$ relations, remain
conjectural.  The coefficient table below completely specifies the new
identity without rounding or an implicitly chosen numerical basis.

\begin{conjecture}[Mixed Gaussian reduction at weight fifteen]
\label{mix:conj:S14}
Let $(X_j,c_j)$ be the ordered rows in Table~\ref{mix:tab:S14}. Then
\begin{equation}\label{mix:eq:S14}
 \sum_{j=0}^{15}c_jX_j=0,
 \qquad\text{equivalently}\qquad
 S_{14}=-\sum_{j=1}^{15}\frac{c_j}{c_0}X_j.
\end{equation}
\end{conjecture}

\begin{table}[htbp]
\centering\small
\begin{tabular}{@{}lr@{}}
\toprule
Period $X_j$ & Integer coefficient $c_j$\\
\midrule
$S_{14}$ & $17137739346472321847995257166233600$ \\
$g_{14,1}$ & $-23326883533329005530721963055513600$ \\
$g_{12,3}$ & $-1241749496845608511840387596288000$ \\
$g_{10,5}$ & $577580508445935074901301710028800$ \\
$g_{8,7}$ & $-584626148672522170960472702976000$ \\
$g_{6,9}$ & $1032524260038200489807505575116800$ \\
$g_{4,11}$ & $-2814967862322082111057500831744000$ \\
$g_{2,13}$ & $16994562809106694872629900633702400$ \\
$\pi^{15}$ & $-969820884575582157963236591$ \\
$\beta(2)\zeta(13)$ & $2074025115468857154352052304000$ \\
$\beta(4)\zeta(11)$ & $23496914452549436838383360716800$ \\
$\beta(6)\zeta(9)$ & $112557931663509122499903946752000$ \\
$\beta(8)\zeta(7)$ & $319433131232914030174904411750400$ \\
$\beta(10)\zeta(5)$ & $765318186561145170729756524544000$ \\
$\beta(12)\zeta(3)$ & $3477679796290052414651630616576000$ \\
$\beta(14)L$ & $34275478692944643695990514332467200$ \\
\bottomrule
\end{tabular}
\caption{The complete primitive integer vector for the conjecture
\eqref{mix:eq:S14}; the first row is $j=0$.}
\label{mix:tab:S14}
\end{table}

For orientation, the first right-side coefficient and the last are
\[
 -\frac{c_1}{c_0}=\frac{28551728212443286}{20976315815914861},
 \qquad
 -\frac{c_{15}}{c_0}=-2.
\]
The intermediate coefficients are part of the displayed vector and
should not be replaced by decimals.  Its entries have greatest common
divisor one.  The coefficient of $\beta(14)L$ matches the value $-2$
in the preceding frozen even-index candidates; this pattern is evidence
for a possible uniform formula, not a proof of one.

\subsection{Discovery and independent numerical audit}

The exploratory evaluation used $3600$ finite Euler weights computed as
exact integers, with $1100$ decimal digits in the subsequent arithmetic.
PSLQ then ran at $850$ digits, tolerance $10^{-800}$, coefficient bound
$10^{65}$, and an iteration cap of $200000$. It returned the primitive
vector in Table~\ref{mix:tab:S14}. After that vector was frozen, its
normalized residual at the retained $1100$-digit precision was
approximately $1.6223\times10^{-1084}$.  This is a numerical diagnostic;
Euler truncation and finite-precision arithmetic are not thereby made
exact. The independent exact certificate below is the authoritative
proximity statement.

A second calculation, at $110$ decimal digits, evaluated the periods
from Mellin integrals rather than the finite Euler sums. It used
\[
 S_{14}=-\frac1{13!}\int_0^\infty
     \frac{t^{13}e^{-t}\log(1+e^{-2t})}{1+e^{-2t}}\,dt
\]
and, for $a\ge2$,
\[
 g_{a,b}=\frac1{(a-1)!}\int_0^\infty t^{a-1}
  \operatorname{Im}\frac{ie^{-t}\operatorname{Li}_b(ie^{-t})}
                                  {1-ie^{-t}}\,dt.
\]
At $b=1$ it used the integrated logarithmic formula
\[
 g_{a,1}=-\frac1{2(a-2)!}\int_0^\infty
 t^{a-2}\log(1+e^{-2t})\arctan(e^{-t})\,dt.
\]
The normalized residual was approximately
$-1.2378\times10^{-110}$. No rigorous quadrature-error enclosure is
claimed for this numerical audit. It checks conventions and gives
independent corroboration of the finite-sum calculation.

\begin{proposition}[Exact rational proximity at weight fifteen]
\label{mix:prop:proximity}
For the actual periods in Table~\ref{mix:tab:S14}, put
$R_{14}=\sum_{j=0}^{15}c_jX_j/c_0$. Then
\begin{equation}\label{mix:eq:proximity}
 \boxed{\qquad |R_{14}|<10^{-775}.\qquad}
\end{equation}
The delivered rational interval contains zero. This proves proximity;
the assertion $R_{14}=0$ remains the conjecture.
\end{proposition}

\subsection{Every analytic premise of the rational certificate}

For a sequence $f_n$, define
\[
 \Delta^jf_0=\sum_{n=0}^j(-1)^n\binom jn f_n,
 \qquad E_N(f)=\sum_{j=0}^{N-1}2^{-j-1}\Delta^jf_0.
\]
A finite interchange of sums and Pascal's identity give
\begin{equation}\label{mix:eq:finite-euler}
 E_N(f)=2^{-N}\sum_{n=0}^{N-1}(-1)^n f_n
                    \sum_{r=n+1}^N\binom Nr.
\end{equation}
The positive weights inside the last sum are exact integers. This
avoids an unstable floating-point finite-difference triangle.

For integer $a\ge2,b\ge1$, the ordinary Mellin formulas give
\[
 \frac{H_{2n}^{(b)}}{(2n+1)^a}
 =\frac1{\Gamma(a)\Gamma(b)}\int_0^1\int_0^1
 \frac{(-\log x)^{a-1}(-\log y)^{b-1}}{1-y}
       \{(x^2)^n-(x^2y^2)^n\}\,dy\,dx.
\]
Keep the two terms grouped, especially at $b=1$. Finite differences
replace them by $(1-x^2)^j-(1-x^2y^2)^j$. Summing the omitted Euler tail
shows that $g_{a,b}-E_N(f)$ is the same positive Mellin measure applied
to
\[
 2^{-N}\{W_N(x^2)-W_N(x^2y^2)\},\qquad
 W_N(u)=\frac{(1-u)^N}{1+u}.
\]
On $[0,1]$, $W_N$ decreases and $|W_N'|\le N+1$. Consequently
\begin{equation}\label{mix:eq:g-bound}
 E_N(f)-\frac{N+1}{2^N3^a}(1+2^{-b})
       \le g_{a,b}\le E_N(f).
\end{equation}
Indeed the difference of $W_N$ values is between
$-(N+1)x^2(1-y^2)$ and zero, and its integral is
$-(N+1)3^{-a}(1+2^{-b})$.
For $f_n=H_n/(2n+1)^p$, replace $y^2$ by $y$, obtaining
\begin{equation}\label{mix:eq:S-bound}
 E_N(f)-\frac{N+1}{2^N3^p}\le S_p\le E_N(f),\qquad p\ge2.
\end{equation}
All interchanges follow either from absolute convergence of the
original series at these integer orders or from the displayed integrable
bound on the grouped Euler kernels.

The single-value sequences $(2n+1)^{-s}$ and $(n+1)^{-s}$ are moments
of positive measures of mass one. Their omitted Euler tails are
nonnegative and at most $2^{-N}$. This encloses $\beta(s)$ and
$\eta(s)$; multiplying by $(1-2^{1-s})^{-1}$ encloses $\zeta(s)$
when $s>1$. Finally, use
\[
 \pi=16\arctan(1/5)-4\arctan(1/239),\qquad
 L=2\sum_{n\ge0}\frac1{(2n+1)3^{2n+1}}.
\]
The first omitted term encloses each alternating arctangent series.
For $M$ terms of the logarithm series, its positive tail is at most
$9/[4(2M+1)3^{2M+1}]$.  Machin's identity uses the real arctangent:
the tangent of $4\arctan(1/5)-\arctan(1/239)$ is one, and this angle
lies in $(0,\pi/2)$, fixing its value as $\pi/4$.

\begin{proof}[Proof of Proposition~\ref{mix:prop:proximity}]
The standalone program \texttt{code/verify\_s14.py} reconstructs all
sixteen period intervals using $N=2600$, $650$ terms in each arctangent
series, and $M=950$ logarithm terms.  For the nested sums it takes the
common denominator $D=\operatorname{lcm}(1,\ldots,2N)$: each harmonic
numerator and each power denominator in \eqref{mix:eq:finite-euler}
can then be accumulated as an integer before one final rational division.
No decimal input, numerical period evaluation, PSLQ calculation or
assumed identity enters the replay.

The verifier reads the frozen integer vector, forms the normalized
interval dot product by exact rational interval arithmetic, and compares
its endpoints to $\pm10^{-775}$. These comparisons are strict, and the
interval contains zero. It also compares every reconstructed basket
interval and both final endpoints against
\texttt{data/s14\_rational\_certificate.json}. This is a finite proof
of the enclosure \eqref{mix:eq:proximity} using only the analytic bounds
just established. It supplies no mechanism for replacing a zero-containing
interval by exact equality.
\end{proof}

\subsection{Provenance, replay and the remaining proof problem}

The common finite-Euler arithmetic is a small adaptation of the incoming
uniform-continuation report's exact proximity implementation, and its
error theorem is retained rather than claimed as new.  The inspected
archive is \texttt{proveit\_polylogarithms\_research\_2026-10-10 (1).zip}
in the pinned repository. Source-member and archive SHA-256 digests are
recorded in \texttt{data/mixed\_gaussian\_provenance.json}.
The new deliverables are the $S_{14}$ vector, its exact proximity
certificate, its independent quadrature audit, and the rational
sampling/finite-part theorems above.

From the package root, the exact certificate is replayed by
\begin{quote}\ttfamily
python code/verify\_s14.py
\end{quote}
It requires only the Python standard library.
Optional diagnostics use \path{code/audit_s14.py},
\path{code/search_s14.py}, and \path{code/verify_mixed_generator.py};
these require \texttt{mpmath}. The search is separate from the
proof of proximity.

A useful next target is a uniform functional identity that implies the
fixed-weight conjectures simultaneously. The existing restricted
single-product obstruction remains valid: it excludes only that stated
row span. The known weight-five octahedral certificate supplies
information outside it. The higher-weight problem is therefore to give
an equally explicit new relation or an independently replayable finite
certificate, not to infer global irreducibility from the restricted
obstruction. The exact resummations above offer consistency constraints
for any proposed uniform formula, while the $S_{14}$ vector supplies a
fresh high-weight target.
